\section{Chương~\ref{sec:eetypes}. Nơ-ron như linh kiện}

Các chế độ hoạt động của từng tế bào đã được đo trực tiếp, bằng dòng điện qua điện cực và ghi điện áp; toán học của bộ tích phân và của việc phân lớp là lý thuyết kinh điển, còn việc bộ não dùng sự tái cấu hình chế độ để tính toán vẫn là giả thuyết.

{\footnotesize
\setlength{\tabcolsep}{3pt}\renewcommand{\arraystretch}{1.15}
\begin{longtable}{>{\raggedright}p{0.5cm}>{\raggedright}p{4.0cm}>{\raggedright}p{5.6cm}>{\raggedright}p{2.3cm}>{\raggedright\arraybackslash}p{1.7cm}}
\toprule
TT & Khẳng định & Thí nghiệm và điều đã đo & Nguồn & Trạng thái \\
\midrule\endfirsthead
\toprule
TT & Khẳng định & Thí nghiệm và điều đã đo & Nguồn & Trạng thái \\
\midrule\endhead
1 & Bộ cộng hưởng: dòng chậm chống lại thay đổi điện áp biến nơ-ron thành bộ lọc thông dải với cực đại ở vài đến vài chục hertz (mục~\ref{sec:resonator}) & Trong lát não, người ta đưa vào nơ-ron \nt{GLUT} dòng hình sin có tần số tăng dần và đo trở kháng $|Z(f)|$: cực đại ở $2$--$5$~Hz tại $33^\circ$C (khoảng $7$~Hz tại $38^\circ$C); cộng hưởng do các dòng kali và dòng cation chậm tạo ra và được dòng natri dai dẳng khuếch đại. Tổng quan cho thấy cùng thủ pháp ở nhiều lớp tế bào & \cite{HuVervaekeStorm2002,HutcheonYarom2000} & đã đo \\
2 & Dòng chậm hoạt động như cuộn cảm, cùng với điện dung màng tạo thành mạch dao động & Đáp ứng dưới ngưỡng của sợi trục với dòng điện đã được đo và mô tả bằng các phương trình độ dẫn tuyến tính hóa; sơ đồ tương đương có một cuộn cảm ``hiện tượng luận'' với giá trị phụ thuộc điện áp & \cite{Mauro1970} & lý thuyết \\
3 & Hằng số thời gian của nhóm có phản hồi $\tau_{\text{hd}}=\tau/(1-w)$~\eqref{eq:perfectint}; với $\tau=0.1$~s và $w=0.995$ là $20$~s & Suy ra trong chương từ phương trình tuyến tính của nhóm; được đề xuất như cơ chế giữ vị trí mắt trong một công trình lý thuyết & suy ra trong chương; \cite{Seung1996} & lý thuyết \\
4 & Bộ tích phân vị trí mắt giữ đầu vào nhờ mạng, chứ không nhờ tính chất của một tế bào riêng lẻ & Ghi nội bào ở cá, trong nhân giữ vị trí mắt: sau một lần xoay nhanh, tần số của nơ-ron thay đổi theo bậc và được giữ; một dòng ngắn vào một tế bào chỉ gây dịch chuyển nhất thời, còn các bậc điện áp vẫn được duy trì cả dưới ngưỡng, nhờ các đầu vào khớp thần kinh & \cite{Aksay2001} & đã đo \\
5 & Bộ tích phân không rò cần được học liên tục để tự điều chỉnh; khi lệch chỉnh, nó hoặc quên, hoặc đi vào bão hòa & Cá được huấn luyện bằng phản hồi thị giác tỉ lệ với $\pm$ vị trí mắt: việc giữ vị trí trở nên không ổn định hoặc bị rò, hằng số thời gian của nơ-ron giảm hơn một bậc, xuống $<1$~s; phản hồi thị giác bình thường dần dần khôi phục độ ổn định & \cite{Major2004} & đã đo \\
6 & Nơ-ron hai trạng thái bền: đầu vào ngắn bật một bình nguyên kéo dài, ức chế tắt nó; tính chất này được bật bởi serotonin (mục~\ref{sec:bistable}) & Ghi nội bào các nơ-ron \nt{ACET} điều khiển cơ ở mèo: một kích thích khớp thần kinh ngắn đưa nơ-ron vào hoạt động kéo dài, một ức chế ngắn đặt lại nó; ở mèo bị cắt tủy sống, trạng thái này xuất hiện sau khi đưa vào tiền chất của serotonin & \cite{Hounsgaard1988} & đã đo \\
7 & Hai lớp: bộ tích phân (tần số tăng từ không) và bộ cộng hưởng (tại ngưỡng lập tức có tần số hữu hạn) & Các lớp được suy ra từ lý thuyết rẽ nhánh. Thí nghiệm: trong lát vỏ não, nơ-ron \nt{GLUT} bắt đầu phát xung với tần số nhỏ tùy ý (loại~1), nơ-ron \nt{GABA} nhanh thì nhảy vọt từ một tần số hữu hạn (loại~2) & \cite{Izhikevich2007,Tateno2004} & đã đo \\
8 & Một nơ-ron là bộ tích phân hoặc bộ phát hiện trùng khớp: ức chế rút ngắn cửa sổ cộng & Trong lát hồi hải mã, cửa sổ trong đó các đầu vào kích thích cộng lại tới ngưỡng ngắn hơn $2$~ms; nó bị rút ngắn bởi ức chế thuận chiều đến ngay sau kích thích; trong đuôi gai, nơi ức chế yếu hơn, cửa sổ rộng hơn & \cite{Pouille2001} & đã đo \\
9 & Đường dây trễ từ sợi trục (bảng~\ref{tab:eetypes}) & Ở cú lợn, đã đo độ trễ trong các sợi trục đi tới bộ phát hiện trùng khớp: chiều dài sợi trục khác nhau cho độ trễ khác nhau, và các bộ phát hiện được điều chỉnh theo hiệu thời gian âm thanh đến & \cite{Carr1990} & đã đo \\
10 & Bộ tạo nhịp khi thức và tế bào im lặng khi ngủ & Nơ-ron \nt{SERO} phát xung gần như đều đặn ban ngày, chậm lại trong giấc ngủ sóng chậm và gần như im lặng trong giấc ngủ REM (chi tiết ở chương~\ref{sec:sleep}) & \cite{JacobsAzmitia1992,McGintyHarper1976} & đã đo \\
11 & Linh kiện do các kênh ion và các kênh quảng bá điều chỉnh chúng quyết định (mệnh đề~\ref{prop:eetypes}) & Đã chứng minh trên các mạng nhỏ: các chất điều biến thay đổi tính chất riêng của nơ-ron và độ mạnh của khớp thần kinh, nên cùng một mạch cho ra các nhịp khác nhau & \cite{Marder2012} & đã đo \\
12 & Bộ não dùng sự tái cấu hình chế độ để tính toán (mệnh đề~\ref{prop:eetypes}) & Chưa có thí nghiệm trực tiếp nào trong đó việc tái cấu hình chế độ của tế bào là cần thiết để giải một nhiệm vụ; chỉ có các thí nghiệm trong đó chất điều biến thay đổi đầu ra của mạch & \cite{Marder2012} & giả thuyết \\
\bottomrule
\end{longtable}}
